\section{\texttt{PHYS\_TO\_VIRT} 与 \texttt{VIRT\_TO\_PHYS} 宏}

在 \texttt{vmm.h} 中定义了两个核心转换宏：

\codefile{src/include/kernel/vmm.h}
\begin{lstlisting}[style=cstyle, caption={地址转换宏（vmm.h）}]
#define KERNEL_VIRT_OFFSET  0xC0000000u

#define PHYS_TO_VIRT(paddr) \
    ((void*)((uint32_t)(paddr) + KERNEL_VIRT_OFFSET))

#define VIRT_TO_PHYS(vaddr) \
    ((uint32_t)(vaddr) - KERNEL_VIRT_OFFSET)
\end{lstlisting}

\begin{whybox}[title={作用域限制}]
这两个宏\textbf{仅适用于内核直接映射区域}。

对于通过 \texttt{vmm\_alloc\_pages()} 动态映射的虚拟地址（如内核堆），
物理页可能是不连续的，$\text{virt} - \hex{C0000000} \neq \text{phys}$。
这些地址需要通过查询页表（\texttt{vmm\_get\_physical()}）来获取物理地址。

简单规则：
\begin{itemize}
  \item PMM 返回的物理地址 → 用 \texttt{PHYS\_TO\_VIRT} 转换后解引用
  \item PDE/PTE 中存储的始终是物理地址（给 CPU/MMU 用）
  \item 内核堆指针（\texttt{kmalloc} 返回值）→ 不能用 \texttt{VIRT\_TO\_PHYS}
\end{itemize}
\end{whybox}

% ============================================================================

